\section{Background and Problem Setup}

\subsection{On-Policy Distillation}

Let $x$ denote an input prompt and $y=(y_1,\ldots,y_T)$ a response sampled autoregressively from a student policy $\student$.
At position $t$, write $y_{<t}=(y_1,\ldots,y_{t-1})$ for the prefix and define the student and teacher next-token distributions as:
\begin{equation}
    p_t(v) \triangleq \student(v\mid x,y_{<t}), \qquad
    q_t(v) \triangleq \teacher(v\mid x,y_{<t}),
    \label{eq:pt-qt}
\end{equation}
over vocabulary tokens $v\in\mathcal{V}$.
OPD computes supervision on trajectories sampled from the current student, usually by minimizing a reverse-KL objective over student-generated trajectories \citep{gu2024minillm,agarwal2024policy,li2026rethinking}:
\begin{equation}
    \mathcal{L}_{\mathrm{OPD}}(\theta)
    =
    \mathbb{E}_{x,\;y\sim\student(\cdot\mid x)}
    \left[
        \sum_{t=1}^{T}
        D_{\mathrm{KL}}\!\left(p_t \,\|\, q_t\right)
    \right].
    \label{eq:opd-token-kl}
\end{equation}
This token-level decomposition drives OPD's sample efficiency: every position provides a dense teacher signal instead of waiting for a final-answer reward.
It also exposes the central assumption: $q_t$ should be a useful local target on the student-visited state $(x,y_{<t})$.
In practice, we approximate this reward on the student's top-$k$ candidate tokens and obtain a per-position reward vector $r_{t,\cdot}$.
Because rewards are already localized by position, a reliability modifier can simply multiply all rewards at position $t$ by a scalar weight before the policy loss.
This is the interface used by \method: reliable prefixes keep their OPD rewards, while drifted prefixes receive attenuated rewards.

\subsection{Dynamic Metrics}

Following the OPD mechanism analysis~\citep{li2026rethinking}, let $S_t^{(p)}=\operatorname{TopK}(p_t,k)$ and $S_t^{(q)}=\operatorname{TopK}(q_t,k)$.
The overlap ratio proposed in that work is
\begin{equation}
    \mathcal{M}_{\mathrm{overlap}}
    =
    \mathbb{E}_{t}\left[
        \frac{|S_t^{(p)}\cap S_t^{(q)}|}{k}
    \right],
    \label{eq:overlap-ratio}
\end{equation}
which measures whether the student and teacher assign high probability to a shared candidate region.
The entropy gap,
\begin{equation}
    \Delta H_t = |H(q_t)-H(p_t)|,
    \label{eq:entropy-gap}
\end{equation}
tracks whether the two models have similar uncertainty on the same visited state.
These metrics are usually reported as diagnostics after training.
\method~turns them into online control signals.
The current primary metric is the per-position overlap ratio, which measures whether student and teacher still share a high-probability candidate region on the same student prefix.
We also evaluate a stricter action-level variant that asks whether the sampled student token is accepted by the teacher's observable top-$p$/top-$k$ region.

\subsection{Long-Horizon Reliability Problem}

The reward at position $t$ depends on $q_t(v)=\teacher(v\mid x,y_{<t})$, where the prefix $y_{<t}$ was produced by the student.
As $t$ grows, the prefix can drift away from trajectories the teacher would naturally generate.
Prior OPD analysis reports suffix-to-prefix instability and decreasing teacher-continuation advantage as the student prefix depth increases \citep{li2026rethinking}.
Thus the same dense reward that helps on moderate traces can become unreliable on long traces. This motivates position-level reliability control.
We estimate whether each student prefix remains teacher-compatible, then allocate dense supervision to reliable prefixes and attenuate rewards after prefix drift.
